=\chapter{激素通信}\label{ch:g8:hormonal-communication}

\omterm{def:g8:puberty:puberty}{青春期}由它们启动
（\cref{def:g8:puberty:hormones}）；\omterm{def:g8:contraception:pill}{避孕药}同样利用了它们
（\cref{def:g8:contraception:pill}）；惊吓时的热浪
在冲刺途中送来过一次
（\cref{ex:g8:nervous-communication:cases}）。这套“邮政
系统”值得单独介绍：负责写信的\omterm{def:g8:hormonal-communication:glands}{腺体}，
运送信件的\omterm{prop:g4:journey-of-food:absorb}{血液}，以及让广播式投递变得精准的寻址手法——
还有
让整套通信保持平衡的\omterm{prop:g8:hormonal-communication:feedback}{反馈}循环。

\section{腺体、激素与靶标}

\begin{definition}[内分泌腺]\label{def:g8:hormonal-communication:glands}
\emph{腺体}\index{腺体}的产物若
\emph{直接释放入血}（而不是像\omterm{def:g4:journey-of-food:tract}{唾液}腺那样经导管排出），
则称为\emph{内分泌腺}，其产物即为\emph{激素}\index{激素}。
主要的“写信人”有：脑底的
\emph{调控腺体}——\omterm{def:g8:puberty:puberty}{青春期}的启动者，也是其他大多数腺体的指挥者；
颈部的\emph{甲状腺}——其激素决定全身的代谢快慢；
\omterm{def:g7:removing-wastes:kidneys}{肾脏}顶部的\emph{肾上腺}——惊吓和冲刺时的警报激素就来自这里；
\emph{胰腺}——其激素调控\omterm{prop:g4:journey-of-food:absorb}{血液}中的\omterm{prop:g7:organs-at-work:bill}{葡萄糖}；
以及
\emph{\omterm{def:g8:reproductive-systems:female}{卵巢}和\omterm{def:g8:reproductive-systems:male}{睾丸}}——
分泌调控循环的激素与发育相关激素，
详见
\cref{ch:g8:reproductive-systems}。
\end{definition}

\begin{proposition}[由受体寻址]\label{prop:g8:hormonal-communication:address}
\omterm{prop:g4:journey-of-food:absorb}{血液}把每种\omterm{def:g8:hormonal-communication:glands}{激素}运送到全身各处
（\cref{ex:g8:puberty:sites}），但每种\omterm{def:g8:hormonal-communication:glands}{激素}只作用于
其\emph{靶器官}\index{靶器官}：即那些带有该信使对应
\emph{受体}（对接位点）的\omterm{def:g5:first-look-at-cells:cell}{细胞}。这是广播式运送、接收端寻址：
同一封信每家都会收到，
只有钥匙吻合的地方才能读懂。
（\omterm{def:g8:nervous-communication:synapse}{突触}的对接机制见
\cref{def:g8:nervous-communication:synapse}，同一原理，
只是尺度在纳米量级。）
\end{proposition}
\admitted

\begin{example}[一个信使的旅程：警报激素]\label{ex:g8:hormonal-communication:alarm}
骑车时险些出事：肾上腺把
警报\omterm{def:g8:hormonal-communication:glands}{激素}释放入血，几秒之内
——\omterm{prop:g4:journey-of-food:absorb}{血液}循环一周
（\cref{ex:g4:heart-and-blood:round}）即可到达——所有装配了相应受体的
器官立即作出反应：心率加快，气道
扩张，\omterm{prop:g7:digestion-and-nutrients:absorption}{肝脏}释放储存的\omterm{prop:g7:organs-at-work:bill}{葡萄糖}
（\cref{ex:g7:digestion-and-nutrients:breadtrace}中的储备被动员），
\omterm{def:g3:bones-and-muscles:muscle}{肌肉}\omterm{def:g4:heart-and-blood:vessels}{血管}扩张，\omterm{def:g4:journey-of-food:tract}{消化道}
收窄，
（\cref{prop:g7:organs-at-work:supply}中的血流重新分配，
由\omterm{def:g8:hormonal-communication:glands}{激素}维持），\omterm{def:g1:the-five-senses:sense}{皮肤}变得苍白、发麻。
一封信，六种回应，一个协调的状态：准备就绪。
\end{example}

\section{靠反馈维持平衡}

\begin{proposition}[反馈原理]\label{prop:g8:hormonal-communication:feedback}
\omterm{def:g8:hormonal-communication:glands}{激素}系统靠
\emph{反馈}\index{反馈}实现自我调节：写信的\omterm{def:g8:hormonal-communication:glands}{腺体}
会\emph{监测自身分泌的结果}，结果达到后便放缓分泌——
像恒温器，而不是水龙头。
血糖就是最好的例子：进餐后，升高的\omterm{prop:g7:organs-at-work:bill}{葡萄糖}本身会触发胰腺的储存\omterm{def:g8:hormonal-communication:glands}{激素}；
\omterm{prop:g7:organs-at-work:bill}{葡萄糖}被储存后，血糖水平下降，分泌也随之减弱；两餐之间，
另一种\omterm{def:g8:hormonal-communication:glands}{激素}则反向释放储存的\omterm{prop:g7:organs-at-work:bill}{葡萄糖}。
血糖水平因此维持在狭窄的范围内，
日夜如此，
无需任何有意识的驱动因素。
\end{proposition}
\admitted

\begin{example}[反馈失灵：糖尿病]\label{ex:g8:hormonal-communication:diabetes}
在\emph{糖尿病}\index{糖尿病}中，储存\omterm{def:g8:hormonal-communication:glands}{激素}的信号中断：胰腺停止分泌，
或靶\omterm{def:g5:first-look-at-cells:cell}{细胞}不再对其作出回应。进餐之后，血糖
持续在\omterm{prop:g4:journey-of-food:absorb}{血液}中堆积；\omterm{def:g7:removing-wastes:kidneys}{肾脏}的重吸收能力不堪负荷，
糖分随之出现在\omterm{def:g7:removing-wastes:kidneys}{尿液}中——
这正是\cref{ex:g7:removing-wastes:thrift}中的经典信号，
如今终于得到解释。治疗的做法是按剂量
分次补充缺失的\omterm{def:g8:hormonal-communication:glands}{激素}，并与进餐时间配合——
相当于用注射器写信，
依靠这一安排，
无数人因此生活得很好。
\end{example}

\begin{remark}[两套系统的协作]\label{rem:g8:hormonal-communication:partnership}
\Cref{prop:g8:nervous-communication:compare}中的电报与
“邮政”实为同一套管理体系。脑底的调控\omterm{def:g8:hormonal-communication:glands}{腺体}
是两者的枢纽：上方是\omterm{def:g5:brain-in-command:system}{神经系统}，下方是\omterm{def:g8:hormonal-communication:glands}{激素}；
骑车者的惊吓先由\omterm{def:g8:nervous-communication:neuron}{神经}传递，再由\omterm{def:g8:hormonal-communication:glands}{激素}持续维持；
\omterm{def:g8:puberty:puberty}{青春期}则在脑的指挥下，依靠\omterm{def:g8:hormonal-communication:glands}{激素}缓慢推进；
\cref{ex:g8:reproduction-and-environment:lambs}中的羔羊
把日照长度（一种\omterm{def:g8:nervous-communication:neuron}{神经}感知）
转化为繁殖\omterm{def:g8:hormonal-communication:glands}{激素}。
快速的“电线”用于速度至关重要之处，
缓慢的“邮政”用于需要持续数月的程序，
两者之间则是转换机构。
\end{remark}

\section{解读激素系统}

\begin{method}[任一激素的四个问题]\label{met:g8:hormonal-communication:read}
遇到任一\omterm{def:g8:hormonal-communication:glands}{激素}时，应确定以下四点：
\begin{enumerate}
  \item \emph{分泌腺}：哪个\omterm{def:g8:hormonal-communication:glands}{腺体}，
        在何种触发\omterm{def:g6:exploring-our-environment:conditions}{条件}下将\omterm{def:g8:hormonal-communication:glands}{激素}释放入血？
  \item \emph{接收者}：哪些\omterm{prop:g8:hormonal-communication:address}{靶器官}带有它的
        受体，收到信号后各自作出什么反应？
  \item \emph{\omterm{prop:g8:hormonal-communication:feedback}{反馈}}：什么结果会被送回，
        以及如何防止过度分泌？
  \item \emph{失效与应用}：信件投递失败时会发生什么，
        医学是否模仿这一机制——
        如\omterm{def:g8:contraception:pill}{避孕药}与\omterm{ex:g8:hormonal-communication:diabetes}{糖尿病}患者的剂量？
\end{enumerate}
\end{method}

\begin{example}[甲状腺节律的方法]\label{ex:g8:hormonal-communication:thyroid}
写信人：甲状腺，受调控\omterm{def:g8:hormonal-communication:glands}{腺体}的指挥。
收信人：几乎每一个\omterm{def:g5:first-look-at-cells:cell}{细胞}——\omterm{def:g8:hormonal-communication:glands}{激素}
决定\cref{prop:g7:what-is-respiration:power}中缓慢“燃烧”的速率。
\omterm{prop:g8:hormonal-communication:feedback}{反馈}：调控\omterm{def:g8:hormonal-communication:glands}{腺体}读取水平，
并调节其“指挥”力度。失效：分泌过少，整台
“发动机”空转——乏力、怕冷，一切都变慢；过量
则“发动机”飞转。应用：缺失的\omterm{def:g8:hormonal-communication:glands}{激素}可每日一片
补充，这是医学上最古老的“邮寄信件”之一。
\end{example}

\section{练习}

\begin{exercise}[$\star$]\label{exo:g8:hormonal-communication:1}
什么使\omterm{def:g8:hormonal-communication:glands}{腺体}成为内分泌腺？
试举四例，并各说明其\omterm{def:g8:hormonal-communication:glands}{激素}的一种作用。
\end{exercise}

\begin{exercise}[$\star$]\label{exo:g8:hormonal-communication:2}
\omterm{prop:g4:journey-of-food:absorb}{血液}输送到各处，\omterm{def:g8:hormonal-communication:glands}{激素}却只在特定处起作用；
试解释这一悖论。
\end{exercise}

\begin{exercise}[$\star$]\label{exo:g8:hormonal-communication:3}
列出警报\omterm{def:g8:hormonal-communication:glands}{激素}六种反应中的四种，
并分别注明各由哪一章介绍。
\end{exercise}

\begin{exercise}[$\star$]\label{exo:g8:hormonal-communication:4}
叙述\omterm{prop:g8:hormonal-communication:feedback}{反馈}原理，并举出其典型实例。
\end{exercise}

\begin{exercise}[$\star$]\label{exo:g8:hormonal-communication:5}
\omterm{ex:g8:hormonal-communication:diabetes}{糖尿病}中哪一环节失效，
由此在\omterm{def:g7:removing-wastes:kidneys}{肾脏}中又出现了什么变化？
\end{exercise}

\begin{exercise}[$\star$]\label{exo:g8:hormonal-communication:6}
\omterm{def:g5:brain-in-command:system}{神经系统}与\omterm{def:g8:hormonal-communication:glands}{激素}系统在何处相接？
举一个接力传递的例子。
\end{exercise}

\begin{exercise}[$\star\star$]\label{exo:g8:hormonal-communication:7}
为什么用恒温器作为\omterm{prop:g8:hormonal-communication:feedback}{反馈}的合适比喻，
而水龙头却不合适？
\end{exercise}

\begin{exercise}[$\star\star$]\label{exo:g8:hormonal-communication:8}
运用\cref{met:g8:hormonal-communication:read}，
分析胰腺的储存\omterm{def:g8:hormonal-communication:glands}{激素}。
\end{exercise}

\begin{exercise}[$\star\star$]\label{exo:g8:hormonal-communication:9}
解释\omterm{ex:g8:hormonal-communication:diabetes}{糖尿病}患者的注射器与\omterm{def:g8:contraception:pill}{避孕药}，
视为同一类行为，依据方法中的第 4 问
作出解释。
\end{exercise}

\begin{exercise}[$\star\star$]\label{exo:g8:hormonal-communication:10}
判断各项由哪一系统主导，并说明理由：寒战的启动；
胡须的生长；血糖的窄幅范围；
膝跳反射。
\end{exercise}

\begin{exercise}[$\star\star$]\label{exo:g8:hormonal-communication:11}
\cref{ex:g8:reproduction-and-environment:lambs}中的绵羊
使光照变化转为羔羊的诞生，
沿\cref{rem:g8:hormonal-communication:partnership}所述的分工，
追溯接力经过的各个职能。
\end{exercise}

\begin{exercise}[$\star\star\star$]\label{exo:g8:hormonal-communication:12}
“受体把广播变成私人线路构成的网络。”
为这句话辩护，然后再加推敲：
\omterm{def:g5:first-look-at-cells:cell}{细胞}要发生何种变化，才能\emph{加入}
\omterm{def:g8:hormonal-communication:glands}{激素}的听众之列——这又说明
\omterm{def:g8:puberty:puberty}{青春期}的\omterm{prop:g8:hormonal-communication:address}{靶器官}是如何学会“倾听”的。
\end{exercise}

\section{问题：血糖日志}

\begin{problem}[{周末问题——一天的血糖，
读作一组书信往来}]\label{pb:g8:hormonal-communication:1}
某连续血糖监测仪的（虚构）一日记录：
夜间平稳；8 点早餐，血糖升高，
随后平缓回落至 10 点附近；12 点运动，血糖短暂下降，
很快得以纠正；漏掉午餐，血糖整个下午
都保持平稳；17 点受惊（骑车险些出事），
血糖急剧而短暂升高；19 点晚餐，
血糖再次升高又回落。

\medskip\noindent\textbf{第一部分——上午。}
\begin{enumerate}
  \item 夜间平稳：哪一组两封信往来
        在主人入睡时守住了血糖的
        稳定范围？
  \item 早餐后的升降：指出触发因素、
        写信人、信件，以及
        负责储存多余血糖的收信人。
  \item 为何下降的\emph{停止}——不会
        跌破稳定范围——这本身是
        \omterm{prop:g8:hormonal-communication:feedback}{反馈}的标志吗？
  \item 10 点的水平与 7 点的一致。
        恒温器这一比喻预测，
        两者之间的分泌过程会怎样？
\end{enumerate}

\medskip\noindent\textbf{第二部分——下午。}
\begin{enumerate}[resume]
  \item 运动后的下降与回升：哪封信
        召唤了\omterm{prop:g7:digestion-and-nutrients:absorption}{肝脏}的“储备”，哪一章
        所述的“存款”正被动用？
  \item 漏掉午餐并未引起明显变化。哪一种
        伙伴\omterm{def:g8:hormonal-communication:glands}{激素}的静默工作
        维系了整个下午？
  \item 受惊后的急剧升高：指出写信人和
        突然增加的糖分的用途——为
        何种情形做准备，据
        \cref{ex:g8:hormonal-communication:alarm}判断？
  \item 为何受惊引起的升高在一小时之内
        消退，而\omterm{ex:g8:hormonal-communication:diabetes}{糖尿病}中的血糖则持续堆积？
\end{enumerate}

\medskip\noindent\textbf{第三部分——门诊。}
\begin{enumerate}[resume]
  \item 同一天的\omterm{ex:g8:hormonal-communication:diabetes}{糖尿病}患者的记录
        每餐都会不同。描述未经治疗时
        早餐后的曲线，以及\omterm{def:g7:removing-wastes:kidneys}{肾脏}的信号
        会出现在何处。
  \item 治疗后的情况：每餐前测得的
        剂量模仿了什么，为何其大小
        须与所吃食物的量相符？
  \item 医生把健康人的记录称为“一场
        你从未听见的对话”。列出
        这一天的发言者——\omterm{def:g8:hormonal-communication:glands}{腺体}、信件、
        接收者——以一张汇总表或
        一段文字的形式呈现。
  \item 用一句话概括本章的寓意：
        哪一类系统能让某个数值
        在一生中保持稳定，而无需驱动因素？
\end{enumerate}
\end{problem}
